\section{Результаты исследования производительности}

В данном разделе представлены результаты тестирования PATRICIA-дерева в сравнении
со стандартным контейнером \texttt{std::map}, а также анализ потребления памяти и
профилирования исправленной программы.

\subsection{Анализ потребления оперативной памяти}

Для проверки ошибок работы с динамической памятью использовалась утилита Valgrind
Memcheck. Запуск выполнялся на наборе из $100\,000$ вставок, $100\,000$ поисков и
$100\,000$ удалений.

\begin{alltt}
HEAP SUMMARY:
    in use at exit: 122,880 bytes in 6 blocks
    total heap usage: 261,459 allocs, 261,453 frees, 9,988,494 bytes allocated

LEAK SUMMARY:
    definitely lost: 0 bytes in 0 blocks
    indirectly lost: 0 bytes in 0 blocks
    possibly lost: 0 bytes in 0 blocks
    still reachable: 122,880 bytes in 6 blocks

ERROR SUMMARY: 0 errors from 0 contexts
\end{alltt}

\begin{table}[!ht]
\centering
\begin{tabular}{|p{5.5cm}|p{4cm}|p{4cm}|}
\hline
\rowcolor{lightgray}
Показатель & PATRICIA & Комментарий \\
\hline
Ошибки памяти & 0 & Некорректных чтений и записей не обнаружено \\
\hline
Утечки пользовательской памяти & 0 bytes & Узлы дерева освобождаются корректно \\
\hline
\texttt{still reachable} & 122\,880 bytes & Буферы стандартной библиотеки ввода-вывода \\
\hline
Пик heap по Massif & 7.335 MB & Набор из $100\,000$ слов \\
\hline
\end{tabular}
\end{table}

Вывод по памяти: утечек в пользовательском коде не найдено. Память, отмеченная как
\texttt{still reachable}, относится к внутренним буферам libstdc++ и не является
ошибкой реализации словаря.

\subsection{Исследование скорости выполнения}

Замеры времени проводились на операциях вставки, поиска и удаления строковых ключей.
Сравнение выполнялось с \texttt{std::map}.

\begin{table}[!ht]
\centering
\begin{tabular}{|r|r|r|r|r|r|r|}
\hline
\rowcolor{lightgray}
& \multicolumn{3}{c|}{PATRICIA, мкс} & \multicolumn{3}{c|}{\texttt{std::map}, мкс} \\
\hline
\rowcolor{lightgray}
$N$ & Insert & Find & Remove & Insert & Find & Remove \\
\hline
1\,000 & 348 & 208 & 465 & 214 & 161 & 303 \\
\hline
10\,000 & 4\,731 & 3\,246 & 8\,264 & 4\,251 & 2\,405 & 3\,711 \\
\hline
100\,000 & 86\,469 & 75\,838 & 100\,506 & 87\,293 & 73\,691 & 81\,991 \\
\hline
500\,000 & 731\,836 & 695\,895 & 764\,450 & 766\,442 & 748\,168 & 785\,446 \\
\hline
1\,000\,000 & 1\,842\,835 & 2\,022\,130 & 2\,038\,648 & 1\,953\,836 & 1\,985\,049 & 2\,322\,171 \\
\hline
\end{tabular}
\end{table}

На малых объёмах \texttt{std::map} быстрее из-за меньших накладных расходов. На больших
наборах PATRICIA начинает выигрывать на части операций: при $N=500\,000$ быстрее поиск
и удаление, при $N=1\,000\,000$ быстрее вставка и удаление.

\subsection{Анализ производительности с помощью gprof}

Для детального анализа времени выполнения была собрана цель \texttt{lab3\_profile}
с флагом \texttt{-pg}. Отчёт \texttt{gprof} для $100\,000$ операций каждого типа
показал следующие основные строки.

\begin{table}[!ht]
\centering
\begin{tabular}{|l|r|r|}
\hline
\rowcolor{lightgray}
Функция & Доля времени, \% & Число вызовов \\
\hline
\texttt{main} & 46.68 & -- \\
\hline
\texttt{TPatriciaTrie::Remove} & 20.01 & 100\,000 \\
\hline
\texttt{TPatriciaTrie::Find} & 17.78 & 100\,000 \\
\hline
\texttt{TPatriciaTrie::Insert} & 15.56 & 100\,000 \\
\hline
\texttt{NormalizeKey} & 0.00 & 400\,000 \\
\hline
\end{tabular}
\end{table}

Профилирование показывает, что основное время тратится на разбор входа и операции
дерева. Удаление является наиболее дорогой операцией PATRICIA, потому что после
физического удаления узла нужно корректно восстановить обратные рёбра.

\subsection{Сравнение исходной и исправленной версии}

\begin{table}[!ht]
\centering
\begin{tabular}{|l|r|r|r|}
\hline
\rowcolor{lightgray}
Версия & Insert, мкс & Find, мкс & Remove, мкс \\
\hline
Исходная & 111\,895 & 111\,360 & 128\,464 \\
\hline
Исправленная & 99\,566 & 109\,004 & 116\,582 \\
\hline
\end{tabular}
\end{table}

Исправления не изменили асимптотику операций дерева. Разница во времени находится
в пределах практических колебаний запусков внутри контейнера; основная польза исправлений
заключается в корректной обработке ошибок сериализации и честном benchmark.

\pagebreak
